\textbf{Phần 1. Trắc nghiệm lựa chọn}
\setcounterpageref{ex}{0}

\begin{ex} %[3D14N72]
Kết quả của phép nhân $30\,000 \times 3$ là
\choice
{$60000$}
{$100000$}
{\True $90000$}
{$80000$}
\loigiai{
$30000 \times 3 = 90000$.\\
\textbf{Đáp án của câu hỏi là C}
}
\end{ex}

\begin{ex} %[3D14N72]
Kết quả của phép chia $100\,000 : 5$ là
\choice
{$15000$}
{\True $20000$}
{$25000$}
{$30000$}
\loigiai{
$100000 : 5 = 20000$.\\
\textbf{Đáp án của câu hỏi là B}
}
\end{ex}

\begin{ex} %[3D14H72]
Kết quả của phép nhân $12\,150 \times 4$ là
\choice
{$48500$}
{$48700$}
{$47600$}
{\True $48600$}
\loigiai{
$12150 \times 4 = 48600$.\\
\textbf{Đáp án của câu hỏi là D}
}
\end{ex}

\begin{ex} %[3D14H72]
Kết quả của phép chia $84\,880 : 4$ là
\choice
{\True $21220$}
{$21210$}
{$21230$}
{$21240$}
\loigiai{
$84880 : 4 = 21220$.\\
\textbf{Đáp án của câu hỏi là A}
}
\end{ex}

\begin{ex} %[3D14H72]
Tính giá trị biểu thức: $25\,000 \times 2 + 15\,000$.
\choice
{$60000$}
{\True $65000$}
{$70000$}
{$55000$}
\loigiai{
$50000 + 15000 = 65000$.\\
\textbf{Đáp án của câu hỏi là B}
}
\end{ex}

\begin{ex} %[3D14H72]
Tính giá trị biểu thức: $(60\,000 - 20\,000) : 4$.
\choice
{\True $10000$}
{$8000$}
{$12000$}
{$15000$}
\loigiai{
$40000 : 4 = 10000$.\\
\textbf{Đáp án của câu hỏi là A}
}
\end{ex}

\begin{ex} %[3D14V72]
Một cửa hàng buổi sáng bán được $14\,200\text{ kg}$ phân bón, buổi chiều bán được số phân bón gấp đôi buổi sáng. Cả ngày bán được số phân bón là
\choice
{$41600\text{ kg}$}
{$43600\text{ kg}$}
{$28400\text{ kg}$}
{\True $42600\text{ kg}$}
\loigiai{
Chiều: $14200 \times 2 = 28400\text{ kg}$.\\
Cả ngày: $14200 + 28400 = 42600\text{ kg}$.\\
\textbf{Đáp án của câu hỏi là D}
}
\end{ex}

\begin{ex} %[3D14V72]
Một khu rừng hình vuông có chu vi $48\,000\text{ m}$. Cạnh của khu rừng là
\choice
{$11000\text{ m}$}
{$13000\text{ m}$}
{\True $12000\text{ m}$ (hay $12\text{ km}$)}
{$14000\text{ m}$}
\loigiai{
$48000 : 4 = 12000\text{ m}$.\\
\textbf{Đáp án của câu hỏi là C}
}
\end{ex}

\begin{ex} %[3D14V72]
Tìm $x$ biết: $x : 5 = 16\,400$.
\choice
{$81000$}
{$83000$}
{$84000$}
{\True $82000$}
\loigiai{
$x = 16400 \times 5 = 82000$.\\
\textbf{Đáp án của câu hỏi là D}
}
\end{ex}

\begin{ex} %[3D14V72]
Tìm $y$ biết: $y \times 3 = 75\,960$.
\choice
{\True $25320$}
{$25310$}
{$25330$}
{$25340$}
\loigiai{
$y = 75960 : 3 = 25320$.\\
\textbf{Đáp án của câu hỏi là A}
}
\end{ex}

\begin{ex} %[3D14V72]
Một cuộn dây cáp dài $64\,000\text{ m}$. Người ta đã dùng $\dfrac{1}{4}$ cuộn dây cáp. Chiều dài dây cáp còn lại là
\choice
{$46000\text{ m}$}
{$50000\text{ m}$}
{\True $48000\text{ m}$}
{$16000\text{ m}$}
\loigiai{
Đã dùng: $64000 : 4 = 16000\text{ m}$.\\
Còn lại: $64000 - 16000 = 48000\text{ m}$.\\
\textbf{Đáp án của câu hỏi là C}
}
\end{ex}

\begin{ex} %[3D14C72]
Hai kho lương thực chứa tất cả $84\,000\text{ kg}$ gạo. Kho thứ nhất chứa $36\,000\text{ kg}$. Hỏi kho thứ hai chứa nhiều hơn kho thứ nhất bao nhiêu ki-lô-gam gạo?
\choice
{$11000\text{ kg}$}
{\True $12000\text{ kg}$}
{$13000\text{ kg}$}
{$14000\text{ kg}$}
\loigiai{
Kho 2: $84000 - 36000 = 48000\text{ kg}$.\\
Kho 2 hơn kho 1: $48000 - 36000 = 12000\text{ kg}$.\\
\textbf{Đáp án của câu hỏi là B}
}
\end{ex}

\textbf{Phần 2. Trắc nghiệm đúng sai}
\setcounterpageref{ex}{0}

\begin{ex} %[3D14H72]
Xét tính đúng sai của các biểu thức tính:
\choiceTF
{\True $20000 \times 4 + 10000 = 90000$.}
{\True $90000 : 3 - 10000 = 20000$.}
{$(40000 + 10000) \times 2 = 80000$.}
{\True $(40000 + 10000) \times 2 = 100000$.}
\loigiai{
\begin{itemize}
    \item $80000 + 10000 = 90000$: Đúng.
    \item $30000 - 10000 = 20000$: Đúng.
    \item $50000 \times 2 = 100000 \neq 80000$. Ý 3 Sai, Ý 4 Đúng.
\end{itemize}
\begin{itemize}
    \item a) \textbf{Đúng}.
    \item b) \textbf{Đúng}.
    \item c) \textbf{Sai}.
    \item d) \textbf{Đúng}.
\end{itemize}
}
\end{ex}

\begin{ex} %[3D14V72]
Một kho có $72\,000\text{ kg}$ phân bón. Lần 1 chuyển $\dfrac{1}{4}$ số phân bón, lần 2 chuyển $25\,000\text{ kg}$. Xét tính đúng sai:
\choiceTF
{\True Lần 1 chuyển được $18\,000\text{ kg}$ phân bón.}
{\True Cả 2 lần chuyển được $43\,000\text{ kg}$ phân bón.}
{Kho còn lại $30\,000\text{ kg}$ phân bón.}
{\True Kho còn lại $29\,000\text{ kg}$ phân bón ($72000 - 43000 = 29000$).}
\loigiai{
\begin{itemize}
    \item Lần 1: $72000 : 4 = 18000\text{ kg}$. Đúng.
    \item Tổng chuyển: $18000 + 25000 = 43000\text{ kg}$. Đúng.
    \item Còn lại: $72000 - 43000 = 29000\text{ kg} \neq 30000$. Ý 3 Sai, Ý 4 Đúng.
\end{itemize}
\begin{itemize}
    \item a) \textbf{Đúng}.
    \item b) \textbf{Đúng}.
    \item c) \textbf{Sai}.
    \item d) \textbf{Đúng}.
\end{itemize}
}
\end{ex}

\textbf{Phần 3. Trả lời ngắn}
\setcounterpageref{ex}{0}

\begin{ex} %[3D14H72]
Tính: $15\,000 \times 4$.
\shortans{60000}
\loigiai{
\textbf{Đáp số:} $60000$\\
$15000 \times 4 = 60000$.
}
\end{ex}

\begin{ex} %[3D14H72]
Tính: $96\,000 : 3$.
\shortans{32000}
\loigiai{
\textbf{Đáp số:} $32000$\\
$96000 : 3 = 32000$.
}
\end{ex}

\begin{ex} %[3D14V72]
Tính giá trị biểu thức: $12\,000 \times 5 - 20\,000$.
\shortans{40000}
\loigiai{
\textbf{Đáp số:} $40000$\\
$60000 - 20000 = 40000$.
}
\end{ex}

\begin{ex} %[3D14V72]
Tính giá trị biểu thức: $84\,000 : 4 + 16\,000$.
\shortans{37000}
\loigiai{
\textbf{Đáp số:} $37000$\\
$21000 + 16000 = 37000$.
}
\end{ex}

\begin{ex} %[3D14V72]
Tìm $x$ biết: $x \times 4 = 64\,800$.
\shortans{16200}
\loigiai{
\textbf{Đáp số:} $16200$\\
$x = 64800 : 4 = 16200$.
}
\end{ex}

\begin{ex} %[3D14C72]
Một khu đất hình chữ nhật có chu vi $60\,000\text{ m}$, chiều dài gấp đôi chiều rộng. Tính chiều dài khu đất theo mét.
\shortans{20000}
\loigiai{
\textbf{Đáp số:} $20000$\\
Nửa chu vi: $60000 : 2 = 30000\text{ m}$.\\
Chiều dài: $30000 : 3 \times 2 = 20000\text{ m}$.
}
\end{ex}
